% ----------------------------------
%   Introduction and structure of the work
% ----------------------------------

\begin{frame}
\frametitle{Introduction}
\begin{scriptsize}
\begin{itemize}
\item An option gives its holder the right, but not the obligation,
to buy (call) or sell (put) an asset at an agreed price $K$ on an agreed date $T$.

\item In 1973 Black, Scholes and Merton priced European options with a PDE
whose solution has a closed form, assuming a \textbf{constant volatility} $\sigma$.

\item After the 1987 crash markets display the \textbf{smile}:
the volatility implied by the prices depends on the strike and the maturity.

\item Dupire (1994) proposed the \textbf{local volatility model},
where the volatility is a deterministic function $\lvol(t,S)$,
which reproduces any arbitrage-free set of prices.

\item To be useful the model must be \textbf{calibrated}:
find the coefficient $\lvol$ of a PDE from finitely many, noisy observations of its solution.
This is an ill-posed \textbf{inverse problem}.
\end{itemize}
\medskip

\textbf{Aim}: build the pricing theory up to the Dupire equation,
and solve its calibration by Tikhonov regularization following

\smallskip
Egger, H., Engl, H. W. (2005). \emph{Tikhonov regularization applied to the inverse problem of option pricing: convergence analysis and rates}. Inverse Problems, 21(3), 1027.

\smallskip
All the code can be found in \url{https://github.com/M315/TFM}.
\end{scriptsize}
\end{frame}

\begin{frame}
\frametitle{Structure of the work}
\begin{scriptsize}
\begin{enumerate}
\item[2--3.] \textbf{Preliminaries.} Probability spaces, filtrations, martingales, Brownian motion,
It\^{o} integral, It\^{o} formula and SDEs.
\item[4.] \textbf{Black-Scholes.} Replication, the BS PDE, its solution, implied volatility and its surface.
\item[5.] \textbf{Arbitrage theory.} PDE and martingale methods, fundamental theorems of asset pricing.
\item[6.] \textbf{Local volatility.} Dupire formula, forward equation, no-arbitrage conditions,
local vs.\ implied volatility.
\item[7.] \textbf{Calibration.} Ill-posedness, Tikhonov functional, adjoint gradient,
convergence rates, discretization and numerical experiments.
\item[A--C.] \textbf{Appendices.} Implied volatility pipeline, calibration details, code.
\end{enumerate}
\end{scriptsize}

\begin{center}
\begin{tikzpicture}[
    node distance=0.35cm,
    box/.style={draw, rounded corners, align=center, font=\scriptsize, minimum height=0.8cm, text width=1.6cm},
    focus/.style={box, fill=blue!15, thick},
    every path/.style={->, >=stealth}]
    \node[box] (sc) {Stochastic calculus};
    \node[box, right=of sc] (bs) {Black-Scholes};
    \node[box, right=of bs] (arb) {Arbitrage};
    \node[focus, right=of arb] (dup) {Dupire PDE};
    \node[focus, right=of dup] (cal) {Calibration};
    \draw (sc) -- (bs); \draw (bs) -- (arb); \draw (arb) -- (dup); \draw (dup) -- (cal);
\end{tikzpicture}
\end{center}
\begin{scriptsize}
This presentation briefly recalls stochastic calculus and Black-Scholes, and focuses on the Dupire PDE and its calibration.
\end{scriptsize}
\end{frame}
